\documentclass[12pt]{article}
\usepackage[utf8]{inputenc}
\usepackage{amsmath, amssymb}

\title{CSE400 – Fundamentals of Probability in Computing}
\date{}

\begin{document}

\maketitle

\section*{Lecture 21 \& 22: Inequalities and Tail Bounds (Fully Expanded Exam Scribe)}

\hrule
\vspace{10pt}

\section*{1. Tail of a Random Variable}

\subsection*{Definition}

The tail of a random variable X is defined as:

\[
P(X > x)
\]

This represents the probability that X takes values larger than x.

\hrule

\subsection*{Why do we care about tails?}

\begin{itemize}
\item Jobs delayed more than 24 hours
\item Packets dropped due to full buffer
\end{itemize}

These correspond to rare but critical events in systems.

\hrule

\subsection*{Key Idea}

\begin{itemize}
\item Mean → average behavior (easy to compute)
\item Tail → rare extreme events (hard to compute)
\end{itemize}

\hrule

\subsection*{Why are tails hard to compute?}

\begin{itemize}
\item Requires summing many small probabilities
\end{itemize}

\hrule

\section*{2. Tail Example 1 (Binomial Case)}

\subsection*{Problem}

Suppose you’re distributing n jobs among n servers at random. What is the probability that a particular server gets ≥ k jobs?

\hrule

\subsection*{Intuition}

\begin{itemize}
\item Jobs are assigned randomly
\item Most servers get few jobs
\item Occasionally, one server gets many jobs
\end{itemize}

This corresponds to a tail event (rare overload).

\hrule

\subsection*{Model}

\[
X \sim \text{Binomial}(n, p)
\]

\hrule

\subsection*{Exact Tail Probability}

\[
P(X \geq k) = \sum_{i=k}^{n} \binom{n}{i} p^i (1 - p)^{n - i}
\]

\hrule

\subsection*{Step-by-Step Logic}

\begin{enumerate}
\item We are interested in the probability that X is at least k.
\item This includes all outcomes where X = k, k+1, ..., n.
\item The probability of exactly i successes is:
\[
\binom{n}{i} p^i (1-p)^{n-i}
\]
\item Therefore, summing over all i ≥ k gives the tail probability.
\end{enumerate}

\hrule

\section*{3. Tail Example 2 (Poisson Case)}

\subsection*{Problem}

Jobs arrive to a datacenter according to a Poisson process with rate λ jobs/hour. What is the probability that ≥ k jobs arrive during the first hour?

\hrule

\subsection*{Intuition}

\begin{itemize}
\item Jobs arrive randomly over time
\item Most hours have moderate arrivals
\item Sometimes, a burst of arrivals occurs
\end{itemize}

This is a tail event (sudden spike in load).

\hrule

\subsection*{Model}

\[
X \sim \text{Poisson}(\lambda)
\]

\hrule

\subsection*{Exact Tail Probability}

\[
P(X \geq k) = \sum_{i=k}^{\infty} e^{-\lambda} \frac{\lambda^i}{i!}
\]

\hrule

\subsection*{Step-by-Step Logic}

\begin{enumerate}
\item Poisson gives probability of exactly i arrivals.
\item Tail event requires summing probabilities from k to infinity.
\end{enumerate}

\hrule

\section*{4. Why Tail Bounds?}

\subsection*{Observation}

Exact tail probabilities are often hard to compute.

\hrule

\subsection*{Exact Computation}

\begin{itemize}
\item Binomial sum
\item Poisson infinite sum
\end{itemize}

These are complicated and long.

\hrule

\subsection*{Idea}

Instead of computing exact probability:

\[
P(X \geq k)
\]

We compute an upper bound:

\[
P(X \geq k) \leq \text{something simple}
\]

\hrule

\subsection*{Key Idea}

Approximate safely instead of computing exactly.

\hrule

\section*{5. Running Example}

\subsection*{Experiment}

Flip a fair coin n times.

\hrule

\subsection*{Distribution}

\[
X = \text{number of heads}
\]
\[
X \sim \text{Binomial}(n, 1/2)
\]

\hrule

\subsection*{Mean}

\[
E[X] = n/2
\]

\hrule

\subsection*{Question}

What is P(X ≥ 3n/4)?

\hrule

\subsection*{Interpretation}

This represents getting unusually many heads, which lies in the tail region.

\hrule

\section*{6. Markov’s Inequality}

\subsection*{Key Idea}

\begin{itemize}
\item Uses only the mean
\item Large values are rare
\end{itemize}

\hrule

\subsection*{Statement}

\[
P(X \geq a) \leq \frac{E[X]}{a}
\]

\hrule

\subsection*{Assumption}

\[
X \geq 0
\]

\hrule

\subsection*{Proof (Detailed Steps)}

\begin{enumerate}
\item Start from expectation:
\[
E[X] = \sum x P(X = x)
\]

\item Consider only terms where x ≥ a: \\
Since x ≥ a,
\[
x P(X=x) \geq a P(X=x)
\]

\item Therefore:
\[
E[X] \geq \sum_{x \geq a} a P(X=x)
\]

\item Factor out a:
\[
E[X] \geq a \sum_{x \geq a} P(X=x)
\]

\item Recognize probability:
\[
\sum_{x \geq a} P(X=x) = P(X \geq a)
\]

\item Hence:
\[
E[X] \geq a P(X \geq a)
\]

\item Rearranging:
\[
P(X \geq a) \leq \frac{E[X]}{a}
\]
\end{enumerate}

\hrule

\subsection*{Example Logic}

If average = 10, then values like 100 are rare.

\hrule

\section*{7. Chebyshev’s Inequality}

\subsection*{Statement}

\[
P(|X - \mu| \geq k) \leq \frac{\text{Var}(X)}{k^2}
\]

\hrule

\subsection*{Proof (Step-by-Step)}

\begin{enumerate}
\item Let Y = (X - μ)²

\item Apply Markov inequality:
\[
P(Y \geq k^2) \leq \frac{E[Y]}{k^2}
\]

\item Substitute:
\[
P((X - \mu)^2 \geq k^2) \leq \frac{E[(X - \mu)^2]}{k^2}
\]

\item Recognize variance:
\[
E[(X - \mu)^2] = \text{Var}(X)
\]

\item Therefore:
\[
P(|X - \mu| \geq k) \leq \frac{\text{Var}(X)}{k^2}
\]
\end{enumerate}

\hrule

\section*{8. Chernoff Bound}

\subsection*{Key Idea}

Use exponential transformation.

\hrule

\subsection*{Steps}

\begin{enumerate}
\item Start with probability:
\[
P(X \geq a)
\]

\item Apply monotonic transformation:
\[
P(e^{tX} \geq e^{ta})
\]

\item Apply Markov inequality:
\[
P(e^{tX} \geq e^{ta}) \leq \frac{E[e^{tX}]}{e^{ta}}
\]

\item Therefore:
\[
P(X \geq a) \leq \frac{E[e^{tX}]}{e^{ta}}
\]
\end{enumerate}

\hrule

\subsection*{Insight}

\begin{itemize}
\item Converts problem into expectation of exponential
\item Leads to tighter bounds
\end{itemize}

\hrule

\section*{9. Poisson Tail Bound}

Used for Poisson distributions using Chernoff method.

\hrule

\section*{10. Jensen’s Inequality}

\subsection*{Convex Function Definition}

\[
f(\lambda x + (1-\lambda)y) \leq \lambda f(x) + (1-\lambda)f(y)
\]

\hrule

\subsection*{Statement}

\[
f(E[X]) \leq E[f(X)]
\]

\hrule

\subsection*{Interpretation}

Function of expectation is less than expectation of function.

\hrule

\section*{11. Summary}

\begin{itemize}
\item Tail = extreme event probability
\item Exact computation = hard
\item Bounds = easier
\end{itemize}

\hrule

\section*{12. Comparison}

\begin{itemize}
\item Markov → uses mean
\item Chebyshev → uses variance
\item Chernoff → strongest bound
\end{itemize}

\hrule

\section*{End of Complete Exam Scribe}

\end{document}